\section{Marco Normativo}

El marco normativo identifica y analiza las disposiciones legales vigentes en el Estado Plurinacional de Bolivia y los estándares técnicos internacionales aplicables al diseño del sistema. Dado que la plataforma gestionará datos de identificación personal de usuarios, registros de activos informáticos y acuerdos comerciales de servicio, resulta imperativo que la solución tecnológica se alinee con los principios constitucionales de privacidad y los estándares de calidad y seguridad de la industria del software.

\subsection{Marco Jurídico Nacional (Bolivia)}

El ordenamiento jurídico boliviano establece las garantías fundamentales para la protección de la información ciudadana y la formalización de las relaciones contractuales y comerciales en entornos digitales:

\begin{itemize}
    \item \textbf{Constitución Política del Estado Plurinacional de Bolivia:} El Artículo 21 (inciso 2) consagra el derecho fundamental de toda persona a la privacidad, intimidad, honra, propia imagen y dignidad. Asimismo, el Artículo 130 establece la Acción de Protección de Privacidad, otorgando a toda persona la facultad de conocer, actualizar o rectificar los datos personales registrados en bancos de datos públicos o privados. El diseño del software incorpora este mandato al restringir la visibilidad de los datos personales del cliente exclusivamente al personal autorizado del taller.
    \item \textbf{Ley N.º 164 (Ley General de Telecomunicaciones y TIC):} Promulgada el 8 de agosto de 2011, esta ley establece el marco regulatorio general para las tecnologías de la información y comunicación. Sus Artículos 6 y 54 reconocen la inviolabilidad y el secreto de las comunicaciones electrónicas, así como la plena validez jurídica y probatoria de los mensajes de datos, documentos electrónicos y firmas digitales. La propuesta cumple esta norma al formalizar la emisión del comprobante digital de servicio y el registro auditable de autorizaciones del cliente.
    \item \textbf{Ley N.º 453 (Ley General de los Derechos de las Usuarias y los Usuarios y de las Consumidoras y los Consumidores):} Promulgada el 4 de diciembre de 2013, esta norma regula las relaciones entre proveedores de servicios y consumidores. En su Artículo 11 (incisos 1, 2 y 3) consagra el derecho de las usuarias y los usuarios a recibir información fidedigna, veraz, oportuna, clara y completa sobre las características, precios y condiciones de los servicios contratados, así como a un trato equitativo y sin discriminación. El sistema propuesto respalda directamente este derecho al ofrecer a los clientes acceso transparente y continuo al estado de su equipo y al detalle pormenorizado de las intervenciones técnicas realizadas.
\end{itemize}

\subsection{Estándares Técnicos Internacionales de Calidad y Seguridad}

En el ámbito de la ingeniería de software, la aplicación de normas internacionales asegura que el artefacto diseñado satisfaga criterios rigurosos de calidad estructural y operativa:

\begin{itemize}
    \item \textbf{Estándar ISO/IEC 25010:2023 (Modelos de Calidad del Producto Software):} Pertenece a la familia SQuaRE (Software Product Quality Requirements and Evaluation) y en su edición de 2023 define un modelo para evaluar la calidad del producto software a través de nueve características principales: adecuación funcional (functional suitability), eficiencia de desempeño (performance efficiency), compatibilidad (compatibility), capacidad de interacción (interaction capability), fiabilidad (reliability), seguridad (security), mantenibilidad (maintainability), flexibilidad (flexibility) y protección frente a riesgos (safety). Para la especificación de requerimientos no funcionales del sistema propuesto, se priorizan la adecuación funcional, la capacidad de interacción, la fiabilidad y la seguridad informacional.
    \item \textbf{Estándares ISO/IEC 27001 e ISO/IEC 27002 (Seguridad de la Información):} La norma ISO/IEC 27001 especifica los requisitos para establecer, implementar, mantener y mejorar un Sistema de Gestión de la Seguridad de la Información (SGSI), mientras que la norma ISO/IEC 27002 actúa como una guía de buenas prácticas que describe controles de seguridad aplicables a la gestión de accesos, protección de contraseñas mediante funciones criptográficas unidireccionales (hashing con algoritmos seguros como Bcrypt), generación de copias de respaldo de la base de datos y políticas de auditoría interna de transacciones.
\end{itemize}

\subsubsection{Matriz Operativa de Evaluación de Calidad del MVP (ISO/IEC 25010:2023)}

Para operacionalizar la evaluación del prototipo funcional (MVP) de acuerdo con el Objetivo Específico 5, la propuesta establece una matriz de métricas e indicadores de calidad fundamentada en el estándar ISO/IEC 25010:2023, la cual se aplicará en el entorno de pruebas durante Taller de Grado II:

\begin{table}[H]
\centering
\caption{Matriz operativa de evaluación de calidad del MVP según ISO/IEC 25010:2023}
\label{tab:matriz-iso25010}
\small
\begin{tabular}{p{2.2cm} p{2.5cm} p{3cm} p{3cm} p{2.3cm}}
\toprule
\textbf{Dimensión de Calidad} & \textbf{Indicador Técnico} & \textbf{Método de Verificación / Cálculo} & \textbf{Instrumento / Entorno} & \textbf{Criterio de Aceptación} \\
\midrule
Adecuación Funcional & Cobertura de Casos de Prueba & (Casos de prueba aprobados / Ejecutados) $\times$ 100 & Script de pruebas funcionales ($\sim$20 flujos núcleo: Recepción, Diagnóstico, Autorización, Reparación, Pruebas, Entrega, Asistente IA, RBAC) & $\geq$ 95\% casos aprobados \\
\addlinespace
Capacidad de Interacción & Puntuación de Usabilidad Percibida & Cuestionario System Usability Scale (SUS) & Escala SUS a muestra piloto (20 participantes: 5 técnicos, 15 clientes) & Puntaje SUS $\geq$ 75/100 \\
\addlinespace
Fiabilidad (Confiabilidad) & Tasa de Fallos Transaccionales & (Transacciones fallidas / Probadas) $\times$ 100 & Logs de aplicación y pruebas de estrés (50 transacciones simuladas) & 0\% fallos en transiciones \\
\addlinespace
Seguridad de la Información & Efectividad de Controles RBAC & (Peticiones no autorizadas bloqueadas / Intentos) $\times$ 100 & Pruebas de penetración y auditoría (30 escenarios: RBAC, UUIDv4, prompt injection) & 100\% intentos bloqueados \\
\addlinespace
Eficiencia de Desempeño (Backend) & Tiempo de Respuesta HTTP Backend & Tiempo promedio de respuesta HTTP en API REST interna & Herramienta ApacheBench / JMeter sobre servidor Laravel & Latencia HTTP $\leq$ 2,0 s \\
\addlinespace
Eficiencia de Desempeño (Asistente IA) & Tiempo de Respuesta API del LLM & Medición del tiempo total de respuesta de la API externa & Logs de peticiones REST a Google Gemini API & Latencia API IA $\leq$ 4,0 s \\
\bottomrule
\end{tabular}
\end{table}